\documentclass[conference]{IEEEtran}
\usepackage{cite}
\usepackage{amsmath,amssymb,amsfonts}
\usepackage{graphicx}
\usepackage{booktabs}
\usepackage{url}
\usepackage{hyperref}

\title{Behavioral Residualization as a Representation for Unsupervised CAN Intrusion Detection}

\author{
\IEEEauthorblockN{Anonymous Authors}
\IEEEauthorblockA{Prepared for arXiv submission}
}

\begin{document}
\maketitle

\begin{abstract}
Controller Area Network (CAN) intrusion detection must handle attacks that reuse legitimate arbitration IDs.
We investigate whether \emph{per-ID behavioral residualization}---z-scoring windowed timing and payload features against per-ID normal baselines---provides a useful representation for unsupervised anomaly detection under such reuse.
Across six detectors (Isolation Forest, One-Class SVM, LOF, HBOS, Elliptic Envelope, Autoencoder) and two corpora (HCRL and ROAD), residual features improve F1 in the majority of detector--task cells relative to raw features under a matched protocol (normal-only training, fixed false-positive-rate thresholding).
Multi-seed evaluation and block-bootstrap confidence intervals support that gains are not single-run artifacts.
We also report deployment-oriented runtime and detection latency, feature ablations, and an explicit coverage boundary: residualization helps targeted legitimate-ID spoofing-style attacks, while novel-ID floods (HCRL DoS) and cross-ID fuzzing remain difficult.
\end{abstract}

\begin{IEEEkeywords}
CAN bus, intrusion detection, anomaly detection, residual representation, automotive security
\end{IEEEkeywords}

\section{Introduction}
In-vehicle networks based on CAN remain a practical target for injection and spoofing attacks.
Many public benchmarks, notably HCRL Car-Hacking, contain presence-based cues (novel IDs, trivial payload constants) that can inflate detector scores without proving behavioral understanding.
ROAD emphasizes verified attacks that reuse legitimate arbitration IDs, better matching stealthy ECU impersonation.

We do \textbf{not} propose a new anomaly algorithm.
Instead, we investigate whether a CAN-specific \emph{representation}---per-ID behavioral residualization---consistently improves multiple unsupervised detectors under a fixed evaluation protocol.

\section{Hypothesis}
\textbf{H1:} Per-ID behavioral residualization improves unsupervised anomaly detection performance across multiple detectors relative to the same detectors on raw window features, under matched splits, calibration, and thresholds.

We pre-register a simple decision rule based on the fraction of detector--task cells with positive residual$-$raw deltas on F1 and ROC-AUC (see experimental supplements).

\section{Method}
\subsection{Window features}
For each arbitration ID we maintain a sliding window of length $w=30$ frames and compute inter-arrival, DLC, and flat payload statistics (14 features; group definitions in code).

\subsection{Behavioral residualization}
On a calibration segment of \emph{normal} windows we estimate per-ID means and standard deviations (minimum sample support; otherwise global fallback).
Residuals are $r_f=(x_f-\mu_{\mathrm{id},f})/(\sigma_{\mathrm{id},f}+\varepsilon)$.
On ROAD, residual statistics are fit \textbf{only on pre-injection} windows.

\subsection{Detectors}
Isolation Forest, One-Class SVM, Local Outlier Factor (novelty), HBOS, Elliptic Envelope, and a shallow autoencoder.
All train on normal-only features; thresholds target approximately 1\% FPR on held-out normals.

\section{Experimental Design}
Datasets: HCRL (DoS, Fuzzy, RPM, gear; 60k-frame prefixes) and ROAD (one capture per attack family; pre-injection protocol).
Factors: representation $\in\{\mathrm{raw},\mathrm{residual}\}$, detector, multi-seed $\{0,\ldots,4\}$.
Metrics: precision, recall, F1, ROC-AUC, PR-AUC, actual FPR, runtime, model size, detection latency.
Statistics: block bootstrap for metric CIs; multi-seed mean$\pm$std; McNemar on paired residual vs raw decisions; paired block-bootstrap on metric deltas.

\section{Results}
\subsection{Multi-detector representation study}
On HCRL, residualization improved F1 in 20/24 detector--task cells and ROC-AUC in 21/24 (Phase~A).
On ROAD, residualization improved F1 in 32/36 cells (Phase~B; multi-seed Phase~C refines uncertainty).
Gains are largest for legitimate-ID spoofing-style tasks (RPM/gear; reverse light; correlated signal).

\subsection{Negative results (retained)}
DoS (novel-ID flood) remains failed for essentially all detectors.
ROAD fuzzing remains weak; residual ROC can decrease.
Some detector--task pairs show residual F1 regressions when raw already performs well (e.g., autoencoder on selected ROAD tasks).

\subsection{Feature importance}
Permutation and leave-one-feature-out analyses (Isolation Forest, residual features) attribute gains primarily to inter-arrival and payload residual dimensions on spoofing tasks (tables/figures in supplement).

\subsection{Deployment metrics}
HBOS and Isolation Forest offer favorable score throughput; deep/kernel methods are heavier.
Detection latency (attack start to first true positive) is reported in milliseconds and frames for residual detectors.

\section{Coverage Boundary}
\begin{itemize}
\item \textbf{Works:} targeted legitimate-AID behavioral deviations (spoofed signals with ID reuse).
\item \textbf{Fails / weak:} novel-ID floods; short cross-ID fuzzing; settings with severe distribution shift relative to pre-injection calibration.
\end{itemize}

\section{Limitations}
No cross-vehicle validation; no CAN-FD; primarily offline calibration; limited hardware diversity; no online residual adaptation in the main results; HCRL presence artifacts; ROAD evaluated with capped frames and one capture per type in primary multi-seed tables unless otherwise noted.

\section{Related Work}
Isolation Forest, one-class SVM, LOF, histogram and covariance methods, and reconstruction models are standard unsupervised detectors.
Residual and baseline-normalized features appear broadly in anomaly detection; our contribution is a CAN-specific per-ID residual representation evaluated multi-detector multi-dataset with explicit failure modes---not a claim of inventing residualization or Isolation Forest.

\section{Conclusion}
Behavioral residualization is a practical CAN representation that consistently improves multiple unsupervised detectors under legitimate-ID reuse, while leaving novel-ID floods and cross-ID fuzzing as open challenges.
We release a reproducible pipeline for arXiv and further workshop/conference revision.

\section*{Reproducibility}
\begin{verbatim}
pip install -e ".[dev]"
python -m experiments.runners.run_phase_a --config experiments/configs/phase_a.yaml
python -m experiments.runners.run_phase_b_road --config experiments/configs/phase_b_road.yaml
python -m experiments.runners.run_phase_c --config experiments/configs/phase_c.yaml
\end{verbatim}
Software: Python 3.9+, scikit-learn, NumPy, pandas, SciPy, PyYAML, matplotlib.
Seeds, configs, and tables are under \texttt{experiments/} and \texttt{tables/}.

\begin{thebibliography}{00}
\bibitem{hcrl} HCRL Car-Hacking Dataset.
\bibitem{road} Verma et al., ROAD: Real ORNL Automotive Dynamometer CAN IDS dataset.
\bibitem{iforest} Liu et al., Isolation Forest.
\end{thebibliography}
\end{document}
